\documentclass[11pt, letterpaper]{article}

% --- Idioma y codificación ---
\usepackage[spanish]{babel}
\usepackage[utf8]{inputenc}
\usepackage[T1]{fontenc}
\usepackage{lmodern}

% --- Matemáticas ---
\usepackage{amsmath, amssymb, amsthm}

% --- Formato ---
\usepackage[margin=2.5cm]{geometry}
\usepackage{parskip}
\usepackage{booktabs}
\usepackage{hyperref}
\hypersetup{colorlinks=true, linkcolor=blue!70!black, urlcolor=blue!70!black}

% --- Entornos de teoremas ---
\theoremstyle{definition}
\newtheorem{definition}{Definición}[section]
\theoremstyle{plain}
\newtheorem{theorem}{Teorema}[section]
\newtheorem{lemma}[theorem]{Lema}
\newtheorem{corollary}[theorem]{Corolario}
\theoremstyle{remark}
\newtheorem{remark}{Observación}[section]

\title{%
  T02 --- Demostración formal:\\[4pt]
  Cota inferior $\Omega(n \log n)$ para ordenamiento por comparación
}
\author{Curso Linux--C--Rust --- B15 Estructuras de datos en C}
\date{}

\begin{document}
\maketitle
\tableofcontents

% ===================================================================
\section{Definiciones}
% ===================================================================

\begin{definition}[Modelo de comparación]
Un algoritmo opera en el modelo de comparación si la única forma de
obtener información sobre los elementos de entrada es mediante
comparaciones de la forma $a_i \leq a_j$ (o $<$, $>$, $=$), cada una
con resultado binario.
\end{definition}

\begin{definition}[Árbol de decisión]
Un árbol de decisión para ordenar $n$ elementos es un árbol binario
donde:
\begin{itemize}
  \item Cada nodo interno contiene una comparación $a_i \leq a_j$.
  \item La rama izquierda corresponde al resultado ``sí''.
  \item La rama derecha corresponde al resultado ``no''.
  \item Cada hoja está etiquetada con una permutación
    $\pi \in S_n$.
\end{itemize}
La ejecución del algoritmo sobre una entrada concreta corresponde a un
camino raíz$\to$hoja en el árbol. La altura del árbol es el número de
comparaciones en el peor caso.
\end{definition}

% ===================================================================
\section{Lema 1: un árbol binario de altura $h$ tiene a lo sumo
  \texorpdfstring{$2^h$}{2\^h} hojas}
% ===================================================================

\begin{lemma}\label{lem:leaves}
Sea $T$ un árbol binario de altura $h$. Entonces $T$ tiene a lo sumo
$2^h$ hojas.
\end{lemma}

\begin{proof}
Por inducción sobre $h$.

\textbf{Base ($h = 0$).} Un árbol de altura~0 es un solo nodo (hoja).
Hojas $= 1 = 2^0$. \checkmark

\textbf{Paso inductivo.} Sea $T$ un árbol de altura $h \geq 1$ con
raíz~$r$. Los subárboles izquierdo y derecho de~$r$ tienen altura a lo
sumo $h - 1$. Por hipótesis inductiva, cada uno tiene a lo sumo
$2^{h-1}$ hojas. Las hojas de~$T$ son exactamente las hojas de sus
subárboles:
\[
  L(T) = L(T_{\text{izq}}) + L(T_{\text{der}})
  \leq 2^{h-1} + 2^{h-1} = 2^h \qedhere
\]
\end{proof}

% ===================================================================
\section{Lema 2: el árbol de decisión tiene al menos $n!$ hojas}
% ===================================================================

\begin{lemma}\label{lem:nfact}
Todo árbol de decisión que ordena correctamente $n$ elementos tiene al
menos $n!$ hojas.
\end{lemma}

\begin{proof}
Por contradicción. Supongamos que el árbol tiene menos de $n!$ hojas.
Como hay $n!$ permutaciones posibles de la entrada, por el principio
del palomar existen dos permutaciones distintas $\pi_1$ y $\pi_2$ que
llegan a la misma hoja~$\ell$.

La hoja~$\ell$ está etiquetada con una única permutación de salida
$\sigma$. Esto significa que el algoritmo produce la misma salida
$\sigma$ para las entradas $\pi_1$ y $\pi_2$.

Pero $\pi_1 \neq \pi_2$ implica que las permutaciones de ordenamiento
correctas son distintas: $\sigma_1 \neq \sigma_2$ (donde $\sigma_i$ es
la permutación que ordena $\pi_i$). El algoritmo aplica $\sigma$ a
ambas, por lo que al menos una queda mal ordenada.

Contradicción con la hipótesis de que el algoritmo ordena
correctamente. Por lo tanto, $L \geq n!$.
\end{proof}

% ===================================================================
\section{Teorema principal:
  \texorpdfstring{$\Omega(n \log n)$}{Omega(n log n)} comparaciones}
% ===================================================================

\begin{theorem}\label{thm:lower-bound}
Todo algoritmo de ordenamiento basado en comparaciones requiere
$\Omega(n \log n)$ comparaciones en el peor caso.
\end{theorem}

\begin{proof}
Sea $T$ el árbol de decisión de un algoritmo que ordena $n$ elementos,
con altura~$h$ ($=$ peor caso de comparaciones).

Por el Lema~\ref{lem:leaves}: $L(T) \leq 2^h$.
Por el Lema~\ref{lem:nfact}: $L(T) \geq n!$.

Combinando:
\[
  n! \leq 2^h
  \quad\Longrightarrow\quad
  h \geq \log_2(n!)
\]

Resta probar que $\log_2(n!) = \Omega(n \log n)$, lo que se hace en la
Sección~\ref{sec:logfact}.
\end{proof}

% ===================================================================
\section{Demostración de
  \texorpdfstring{$\log_2(n!) = \Theta(n \log n)$}{log2(n!) = Theta(n log n)}}
\label{sec:logfact}
% ===================================================================

\subsection{Método 1: cota inferior directa (sin Stirling)}

\begin{lemma}\label{lem:logfact-lower}
$\log_2(n!) \geq \dfrac{n}{2}\,\log_2\!\left(\dfrac{n}{2}\right)$.
\end{lemma}

\begin{proof}
En el producto $n! = 1 \cdot 2 \cdot 3 \cdots n$, los factores
$\lceil n/2 \rceil, \lceil n/2 \rceil + 1, \ldots, n$ son todos
$\geq n/2$. Hay al menos $\lfloor n/2 \rfloor$ tales factores.
Descartando los factores menores:
\[
  n! = \prod_{i=1}^{n} i
  \geq \prod_{i=\lceil n/2 \rceil}^{n} i
  \geq \left(\frac{n}{2}\right)^{\!\lfloor n/2 \rfloor}
  \geq \left(\frac{n}{2}\right)^{\!n/2}
\]

Tomando $\log_2$:
\[
  \log_2(n!) \geq \frac{n}{2}\,\log_2\!\left(\frac{n}{2}\right)
  = \frac{n}{2}(\log_2 n - 1)
\]

Para $n \geq 4$: $\log_2 n \geq 2$, así que
$\log_2 n - 1 \geq \frac{1}{2}\log_2 n$:
\[
  \log_2(n!) \geq \frac{n}{2} \cdot \frac{1}{2}\log_2 n
  = \frac{n \log_2 n}{4}
  \quad\Longrightarrow\quad
  \log_2(n!) = \Omega(n \log n) \qedhere
\]
\end{proof}

\subsection{Método 2: cota superior (para obtener $\Theta$)}

\begin{lemma}\label{lem:logfact-upper}
$\log_2(n!) \leq n \log_2 n$.
\end{lemma}

\begin{proof}
Cada factor del producto es $\leq n$:
\[
  n! = \prod_{i=1}^{n} i \leq n^n
  \quad\Longrightarrow\quad
  \log_2(n!) \leq \log_2(n^n) = n \log_2 n \qedhere
\]
\end{proof}

Combinando ambas cotas:
\[
  \frac{n \log_2 n}{4} \leq \log_2(n!) \leq n \log_2 n
  \quad\Longrightarrow\quad
  \log_2(n!) = \Theta(n \log n)
\]

\subsection{Método 3: aproximación de Stirling (cota más precisa)}

\begin{theorem}[Stirling]
Para $n \geq 1$:
\[
  n! = \sqrt{2\pi n}\left(\frac{n}{e}\right)^{\!n}
  \left(1 + O\!\left(\frac{1}{n}\right)\right)
\]
\end{theorem}

Derivación de $\log_2(n!)$ usando Stirling:
\[
  \log_2(n!)
  = \log_2\!\left(\sqrt{2\pi n}\right)
  + n\,\log_2\!\left(\frac{n}{e}\right)
  + O\!\left(\frac{1}{n}\right)
\]

Expandimos cada término:
\begin{align*}
  &= \frac{1}{2}\log_2(2\pi n) + n\log_2 n - n\log_2 e
    + O\!\left(\frac{1}{n}\right) \\
  &= n\log_2 n - n\log_2 e + \frac{1}{2}\log_2(2\pi)
    + \frac{1}{2}\log_2 n + O\!\left(\frac{1}{n}\right)
\end{align*}

Como $\log_2 e \approx 1.4427$:
\[
  \log_2(n!) = n\log_2 n - 1.4427\,n + O(\log n)
\]

\begin{table}[h]
\centering
\caption{Verificación numérica de la aproximación de Stirling.}
\begin{tabular}{ccc}
\toprule
$n$ & $\log_2(n!)$ exacto & $n\log_2 n - 1.443n$ \\
\midrule
10   & 21.79   & 18.78  \\
100  & 524.76  & 520.57 \\
1000 & 8529.66 & 8527.14 \\
\bottomrule
\end{tabular}
\end{table}

Conclusión:
$\log_2(n!) = n\log_2 n - \Theta(n) = \Theta(n \log n)$.

% ===================================================================
\section{Corolario: optimalidad de merge sort y heapsort}
% ===================================================================

\begin{corollary}\label{cor:optimal}
Los algoritmos merge sort y heapsort son asintóticamente óptimos en el
modelo de comparación.
\end{corollary}

\begin{proof}
Ambos ejecutan $O(n \log n)$ comparaciones en el peor caso. La cota
inferior demuestra que ningún algoritmo basado en comparaciones puede
hacer menos de $\Omega(n \log n)$. Por lo tanto, su complejidad de
$\Theta(n \log n)$ es óptima.
\end{proof}

\begin{remark}
Quicksort tiene peor caso $O(n^2)$, así que no es óptimo en el peor
caso. Sin embargo, su caso promedio es $O(n \log n)$, que coincide con
la cota inferior. Quicksort randomizado tiene esperanza $O(n \log n)$,
que es óptimo en expectativa.
\end{remark}

% ===================================================================
\section{Corolario: los sorting lineales no contradicen la cota}
% ===================================================================

\begin{corollary}\label{cor:linear}
Counting sort ($O(n + k)$), radix sort ($O(d(n + k))$) y bucket sort
($O(n)$ esperado) no contradicen la cota de $\Omega(n \log n)$.
\end{corollary}

\begin{proof}
Estos algoritmos \textbf{no} operan en el modelo de comparación:
\begin{itemize}
  \item Counting sort usa valores como índices de array.
  \item Radix sort examina dígitos individuales.
  \item Bucket sort distribuye por rangos.
\end{itemize}
La cota $\Omega(n \log n)$ solo aplica a algoritmos que obtienen
información exclusivamente mediante comparaciones binarias. Los
algoritmos lineales explotan propiedades adicionales de los datos
(rango acotado, representación digital), lo que les permite esquivar
la cota.
\end{proof}

% ===================================================================
\section{Resumen de resultados}
% ===================================================================

\begin{table}[h]
\centering
\begin{tabular}{llp{5.5cm}}
\toprule
\textbf{Resultado} & \textbf{Técnica} & \textbf{Clave} \\
\midrule
Árbol de altura~$h$: $\leq 2^h$ hojas
  & Inducción
  & $L(T_{\text{izq}}) + L(T_{\text{der}}) \leq 2^{h-1} + 2^{h-1}$ \\
Árbol de decisión: $\geq n!$ hojas
  & Principio del palomar
  & Dos permutaciones distintas $\to$ misma hoja $\to$ error \\
$h \geq \log_2(n!)$
  & Combinar los dos lemas
  & $n! \leq 2^h$ \\
$\log_2(n!) = \Omega(n \log n)$
  & Mitad de factores $\geq n/2$
  & $(n/2)^{n/2} \leq n!$ \\
$\log_2(n!) = n\log_2 n - 1.443n + O(\log n)$
  & Stirling
  & Asíntota precisa \\
Merge sort / heapsort óptimos
  & Cota inferior vs superior
  & $\Theta(n \log n)$ ambas \\
\bottomrule
\end{tabular}
\end{table}

\end{document}
